% ============================================================
%  State vector layout
% ============================================================

The Kalman filter operates on a state vector $\stateV_k \in \mathbb{R}^{n_k}$
whose dimension $n_k$ \emph{varies across epochs}: per-satellite and
per-(satellite, signal) substates are added when a new sat/signal is
first observed and pruned (or marked stale) when the satellite drops
below the elevation mask or its phase track is interrupted by a cycle
slip without recovery.  All bookkeeping --- key$\to$index maps, free-list
re-use, and dimensional reshaping of $\stateV$ and $\covM$ --- is
encapsulated in \texttt{StateManager}; \texttt{KalmanFilter} consumes
the state through that interface and never indexes raw memory by
hard-coded offsets.

\subsection{Schematic layout}

In schematic form, suppressing the epoch index $k$:
\begin{align}
\stateV \;=\; \big[\,
& \underbrace{\rxpos_\rx^{\Tr}}_{3} \,,\,
  \underbrace{\bm{v}_\rx^{\Tr}}_{0\text{ or }3} \,,\,
  \underbrace{\sLight\,\mathrm{d}t_\rx}_{1} \,,\,
  \underbrace{(\sLight\,\delta t^{(g)}_\rx)_{g\neq g_{\mathrm{ref}}}^{\Tr}}_{|\sysset|-1} \,,\, \notag\\
& \underbrace{\zwd_\rx,\,(g_N, g_E)}_{1+0\text{ or }1+2} \,,\,
  \underbrace{(\sLight\,\mathrm{d}t^{\phi}_\rx)}_{0\text{ or }1} \,,\, \notag\\
& \underbrace{(I^\sat_{\rx,1})_{\sat \in \mathcal{V}_k}^{\Tr}}_{|\mathcal{V}_k|} \,,\,
  \underbrace{(N^\sat_{\rx,\sigl})_{(\sat,\sigl) \in \mathcal{A}_k}^{\Tr}}_{|\mathcal{A}_k|} \,,\, \notag\\
& \underbrace{(B_{\rx,\sigl})_{\sigl \in \sigset}^{\Tr}}_{|\sigset|} \,,\,
  \underbrace{(B^\sat_\sigl)_{(\sat,\sigl) \in \mathcal{D}_k}^{\Tr}}_{|\mathcal{D}_k|}
\,\big]^{\Tr}.
\end{align}
where
\begin{itemize}[leftmargin=*, itemsep=0pt]
\item $\mathcal{V}_k$ is the set of satellites currently visible at
      epoch $k$ (i.e.\ above the elevation mask, with valid orbit/clock
      products and at least one tracked signal),
\item $\mathcal{A}_k \subseteq \mathcal{V}_k \times \sigset$ is the set
      of phase-tracked (satellite, signal) pairs,
\item $\mathcal{D}_k \subseteq \mathcal{V}_k \times \sigset$ is the
      subset of (satellite, signal) pairs for which the BIA product
      lacks a satellite OSB and a state-space estimate is therefore
      created (see \texttt{estimation.estimate\_sat\_biases}).
\end{itemize}
For brevity the schematic suppresses three conditional refinements
detailed below: the rate states that accompany velocity (a receiver
clock-rate and one inter-system clock-rate offset per non-reference
constellation), the nine-state inertial augmentation of \Cref{sec:imu},
and the fact that the entire receiver-side clock and inter-system-offset
block is removed under the BSSD parameterisation (\Cref{sec:bssd}).

\subsection{Component-by-component description}

\paragraph{Position $\rxpos_\rx \in \mathbb{R}^3$.}
ECEF Cartesian receiver position.  Always present.  Initialised either
from the RINEX header approximate position or, after the first epoch,
from a Weighted Least-Squares (WLS) seed solution to within the
linearisation radius of the geometric range model
(\Cref{sec:obs-geometry}).

\paragraph{Velocity $\bm{v}_\rx \in \mathbb{R}^3$.}
Optional.  Enabled by \texttt{estimation.estimate\_velocity = true}, or
forced on by tightly-coupled IMU integration.  When present, position
and velocity form a coupled continuous-white-noise-acceleration (CWNA,
integrated-random-walk) block driven by a \emph{single} acceleration
spectral density \texttt{process\_noise.acceleration\_psd}
(\Cref{eq:kf-Q-cwna}); the scalar \texttt{position\_q} of the
velocity-free case is then unused.  Enabling velocity also introduces a
receiver clock-rate state $\sLight\,\dot{\mathrm{d}t}_\rx$, and one
inter-system clock-rate offset per non-reference constellation, each
CWNA-coupled to its bias counterpart (\texttt{clock\_rate\_q},
\texttt{system\_time\_rate\_offset\_q}).

\paragraph{Receiver clock $\sLight\,\mathrm{d}t_\rx$.}
Single state, in metres.  Process noise is large
(\texttt{clock\_bias\_q} $\sim 10^3\,\text{m}^2$ per second) so the
state is effectively re-estimated each epoch; this captures
oscillator drift without imposing a stochastic model that may not
match the receiver's actual oscillator.  Under the BSSD
(between-satellite single-difference) parameterisation
(\Cref{sec:bssd}) the receiver clock is unobservable and the state is
dropped entirely --- as are the inter-system offsets below --- so the
whole receiver-side clock block vanishes in PPP-RTK SD mode.

\paragraph{Inter-system time offsets $\sLight\,\delta t^{(g)}_\rx$.}
One per active constellation other than the chosen reference (default
GPS).  These absorb the per-constellation clock biases that arise from
mixing analysis-centre-defined clock products across systems
(``GAL\,$-$\,GPS time offset,'' etc.).

\paragraph{Tropospheric ZWD and gradients.}
$\zwd_\rx$ models the zenith wet delay; the zenith hydrostatic
component $\zhd_\rx$ is treated deterministically via the Saastamoinen
model \cite{saastamoinen1972atmospheric}.  Total slant tropospheric
delay is then
\begin{equation}
T^\sat_\rx \;=\;
m_h(\theta^\sat_\rx)\,\zhd_\rx
\;+\;
m_w(\theta^\sat_\rx)\,\zwd_\rx
\;+\;
m_w(\theta^\sat_\rx)\cot\theta^\sat_\rx \big[
g_N \cos\alpha^\sat_\rx + g_E \sin\alpha^\sat_\rx
\big],
\label{eq:tropo-mapping}
\end{equation}
with elevation $\theta^\sat_\rx$, azimuth $\alpha^\sat_\rx$, and the
Niell or VMF mapping functions $m_h, m_w$.  Gradients $(g_N, g_E)$ are
optional.

\paragraph{Per-satellite slant ionosphere $I^\sat_{\rx,1}$.}
PPP-RTK-only.  One state per visible satellite, expressed at the L1
frequency; signal-specific delays follow \Cref{eq:iono-scaling}.
Initial value seeded from the Klobuchar broadcast model when no
external ionospheric product is available.  Its process noise is a
random walk whose zenith spectral density \texttt{ionosphere\_q} is
projected onto the slant by the squared obliquity factor
$1/\sin^2\theta^\sat_\rx$ (\Cref{eq:kf-Q-iono}), letting low-elevation
states vary faster.  Modelling these per-sat states explicitly is what
permits PPP-RTK to retain integer ambiguity recovery without forming
the ionosphere-free combination (which amplifies code noise by a factor
of $\sim 3$).

\paragraph{Carrier-phase ambiguities $N^\sat_{\rx,\sigl}$.}
One state per phase-tracked $(\sat, \sigl)$ pair, in cycles.  Treated
as float during the prediction/update cycle and resolved to integers
post-hoc by the LAMBDA-family algorithms in
\texttt{filter/AmbiguityResolver} \cite{teunissen1995lambda} when AR
is enabled.  Cycle slips trigger a re-seed via
\texttt{StateManager::resetAmbiguity}, which inflates the diagonal
covariance entry while preserving the cross-state correlations that
permit a consistent re-fix.

\paragraph{Receiver code biases $B_{\rx,\sigl}$.}
These absorb the receiver-side inter-frequency biases that the OSB
convention does not reference to zero.  They are \emph{not} one per
signal: \texttt{StateInitializer::resolveReceiverBias} keys them per
(system, frequency band) as inter-frequency biases
(\texttt{"<system>\_IFB<band>"}) plus, in the alternative scheme, a
single per-system differential code bias (\texttt{"<system>\_DCB"}).
Like the receiver clock they are suppressed under BSSD.

\paragraph{Satellite code biases $B^\sat_\sigl$.}
Per (satellite, signal) pair, but \emph{only when the BIA file does
not supply an OSB} for that pair --- typically a third-frequency
signal in a triple-frequency setup.  When present, these states are
weakly stochastic (\texttt{sat\_bias\_q} $\sim 10^{-10}$) since the
underlying physical biases are nearly stationary on the daily
timescale.

\paragraph{Inertial augmentation.}
When a tightly-coupled IMU is active the core state grows by nine
further entries --- a three-axis attitude error and the accelerometer
and gyroscope biases --- propagated by the strapdown mechanisation
rather than the random-walk process model.  These are treated in full
in \Cref{sec:imu}.

\paragraph{Baseline differential troposphere.}
In baseline mode the per-station tropospheric, clock, ionosphere, and
bias states are absent; an optional single differential zenith-wet-delay
state $\Delta\zwd_{\rx_{\mathrm{rov}}\rx_{\mathrm{ref}}}$ may be carried
(\texttt{baseline.estimate\_differential\_tropo}, \Cref{sec:dd}).

\subsection{Dynamic indexing}\label{sec:state-indexing}

The variable-dimension blocks are addressed via key$\to$index maps
maintained by \texttt{StateManager}.  The canonical key formats are
defined once in \texttt{include/filter/StateKeySchema.hpp}, each with a
\texttt{make}/\texttt{parse} pair:
\begin{itemize}[leftmargin=*, itemsep=0pt]
\item iono key: \texttt{"<sys><prn>"}, e.g.\ \texttt{"G05"} (delimiter-free);
\item ambiguity key: \texttt{"<sys><prn>\_<signal>"}, e.g.\ \texttt{"G05\_L1C"}
      --- split on the \emph{first} underscore, so an iono-free combination
      \texttt{"G05\_L1C\_L5Q"} retains its compound signal;
\item satellite signal-bias key: \texttt{"<sys><prn>|<signal>"}, e.g.\
      \texttt{"G05|C1W"} --- a \emph{pipe} delimiter and a distinct map,
      \emph{not} the underscore ambiguity namespace;
\item receiver code-bias key: \texttt{"<system>\_IFB<band>"} or
      \texttt{"<system>\_DCB"};
\item inter-system time-offset key: the bare system name
      (\texttt{"Galileo"}, \dots; GPS is the reference and never appears);
\item double-difference ambiguity key (baseline):
      \texttt{"DD|<cons>|<sat>|<code>|<pivot>"}.
\end{itemize}
This single source of truth eliminates the off-by-one errors that
plague hand-managed dynamic state vectors and makes the BSSD/DD
S-transforms (\Cref{sec:diff}) implementable as block operations on the
relevant key sets.

\paragraph{Ambiguity basis.}
By default the ambiguity block carries one per-band float ambiguity per
tracked (satellite, signal) in metres.  Setting
\texttt{estimation.state\_ambiguity\_basis = "nl\_wl"} switches the
whole run to a narrow-lane\,/\,wide-lane parameterisation: per satellite
a narrow-lane state (key \texttt{"<sat>\_NL\_<band>"}) and one wide-lane
state per band pair (\texttt{"<sat>\_WL\_<a>\_<b>"}, semantically
$N_a - N_b$), both in \emph{cycles} rather than metres so the wide-lane
and \textsc{lambda} arithmetic stays native to integers, with the
wavelength scaling moved into the Jacobian.  The two bases never coexist
in a single state vector.
